% English translation of questions/5780/semA-moedB/q1.tex (metadata is taken from the Hebrew file).

\begin{question}
Investigate the function $y=\frac{x^3-1}{4x^2}$ according to the following items:
\begin{enumerate}
  \item Domain, intersection points with the axes, evenness, oddness
  \item Continuity, points of discontinuity
  \item Asymptotes
  \item Intervals of increase and decrease, extremum points
  \item Intervals of convexity and concavity, inflection points
  \item Graphical description.
\end{enumerate}
\end{question}

\begin{hint}
It is convenient to write the function in the form $y=\frac{x}{4}-\frac{1}{4x^2}$: from it one immediately sees the oblique asymptote, and differentiation is simple.
\end{hint}

\begin{solution}
Denote $f(x)=\frac{x^3-1}{4x^2}=\frac{x}{4}-\frac{1}{4x^2}$.
\begin{enumerate}
  \item \textbf{Domain:} the denominator vanishes only at $x=0$, hence $D=\{x\in\R : x\neq 0\}$.

  \textbf{Intersection with the axes:} $x=0$ is not in the domain, hence there is no intersection with the $y$-axis. Intersection with the $x$-axis: $x^3-1=0 \iff x=1$, i.e. the point $(1,0)$.

  \textbf{Evenness:} $f(-x)=\frac{-x^3-1}{4x^2}$. For example $f(1)=0$ whereas $f(-1)=-\frac12$, so that $f(-1)\neq f(1)$ and also $f(-1)\neq -f(1)$. Hence the function is neither even nor odd.

  \item $f$ is a quotient of polynomials, hence continuous on its entire domain. The only point where $f$ is not defined is $x=0$. There the numerator tends to $-1$ and the denominator $4x^2\to 0^+$, hence
  \[ \lim_{x\to 0^-} f(x)=\lim_{x\to 0^+} f(x)=-\infty , \]
  i.e. at $x=0$ there is a discontinuity of the second kind (infinite).

  \item \textbf{Vertical asymptote:} by the previous item, $x=0$ is a vertical asymptote (from both sides the function tends to $-\infty$). There are no other vertical asymptotes since $f$ is continuous at every other point.

  \textbf{Oblique asymptote:}
  \[ m=\lim_{x\to\pm\infty}\frac{f(x)}{x}=\lim_{x\to\pm\infty}\frac{x^3-1}{4x^3}=\frac14,\qquad
     n=\lim_{x\to\pm\infty}\Big(f(x)-\frac{x}{4}\Big)=\lim_{x\to\pm\infty}\Big(-\frac{1}{4x^2}\Big)=0 . \]
  Hence $y=\frac{x}{4}$ is an oblique asymptote both at $+\infty$ and at $-\infty$. Since $f(x)-\frac{x}{4}=-\frac{1}{4x^2}<0$, the graph lies below the asymptote. There is no horizontal asymptote (since $f(x)\to\pm\infty$ as $x\to\pm\infty$).

  \item Differentiating:
  \[ f'(x)=\frac14+\frac{1}{2x^3}=\frac{x^3+2}{4x^3}. \]
  $f'(x)=0 \iff x^3=-2 \iff x=-\sqrt[3]{2}$. Sign table (the numerator $x^3+2$ is positive for $x>-\sqrt[3]{2}$, the denominator $4x^3$ is positive for $x>0$):
  \begin{itemize}
    \item $x<-\sqrt[3]{2}$: numerator negative, denominator negative, $f'>0$ --- $f$ is increasing.
    \item $-\sqrt[3]{2}<x<0$: numerator positive, denominator negative, $f'<0$ --- $f$ is decreasing.
    \item $x>0$: numerator positive, denominator positive, $f'>0$ --- $f$ is increasing.
  \end{itemize}
  Hence $f$ is increasing on $(-\infty,-\sqrt[3]{2})$ and on $(0,\infty)$, and decreasing on $(-\sqrt[3]{2},0)$.
  At $x=-\sqrt[3]{2}$ the derivative changes from positive to negative, hence this is a local maximum point:
  \[ f(-\sqrt[3]{2})=\frac{-2-1}{4\sqrt[3]{4}}=-\frac{3}{4\sqrt[3]{4}}\approx -0.47 . \]
  There are no other extremum points (at $x=0$ the function is not defined).

  \item
  \[ f''(x)=\Big(\frac14+\frac12x^{-3}\Big)'=-\frac{3}{2x^4}. \]
  For every $x\neq 0$ we have $f''(x)<0$. Hence $f$ is concave ($\cap$) on each of the intervals $(-\infty,0)$ and $(0,\infty)$, and is not convex ($\cup$) on any interval. The sign of $f''$ does not change within the domain, hence there are no inflection points.

  \item \textbf{Graphical description:} on the left side the graph lies below the asymptote $y=\frac{x}{4}$ and hugs it as $x\to-\infty$. It increases up to the maximum $\big(-\sqrt[3]{2},-\frac{3}{4\sqrt[3]{4}}\big)\approx(-1.26,-0.47)$, and then decreases to $-\infty$ as $x\to0^-$. On the right side it increases from $-\infty$ (as $x\to0^+$), crosses the $x$-axis at $(1,0)$, and approaches the asymptote $y=\frac{x}{4}$ from below as $x\to\infty$. The graph is concave everywhere.
\end{enumerate}
\end{solution}
